\documentclass[conference]{IEEEtran}

\usepackage{graphicx}
\usepackage{amsmath,amssymb}
\usepackage{float}
\usepackage{booktabs}
\usepackage{multirow}

\begin{document}

\title{
    \textbf{Benchmark Report: 508.namd\_r} \\
    \text{Register Spill Analysis on RISC-V via gem5}
}

\author{
    \IEEEauthorblockN{Lütfullah Burak Kaya}
    \IEEEauthorblockA{
        Ozyegin University \\
        burak.kaya.50711@ozu.edu.tr
    }
    \and
    \IEEEauthorblockN{Ismail Akturk}
    \IEEEauthorblockA{
        Ozyegin University \\
        ismail.akturk@ozyegin.edu.tr
    }
}

\newcommand{\LongDate}{%
    \number\day\space
    \ifcase\month
    \or January\or February\or March\or April\or May\or June%
    \or July\or August\or September\or October\or November\or December%
    \fi
    \space \number\year
}

\twocolumn[
\begin{@twocolumnfalse}
\maketitle
\begin{center}{\small \LongDate}\end{center}
\vspace{1em}
\end{@twocolumnfalse}
]

% ================================================

\begin{abstract}
This report presents the register spill characterization of the
\texttt{508.namd\_r} benchmark from SPEC CPU2017, executed on a RISC-V
(RV64GC) target using the gem5 TimingSimpleCPU simulator.
A custom SpillDetector was integrated into gem5 to count store-then-load
patterns on the stack within a user-defined Region of Interest (ROI).
The test dataset (1 iteration) was used to obtain a fast yet representative
measurement. Results show a spill rate of approximately 1.03\% within the
ROI, corresponding to over 327 million detected spill events out of
31.9 billion simulated instructions.
\end{abstract}

% =========================================================
\section{Benchmark Overview}

\texttt{508.namd\_r} is the SPEC CPU2017 rate variant of NAMD, a molecular
dynamics application used to simulate large biomolecular systems.
The benchmark models the ApoA1 lipid-protein complex and is
compute-intensive, dominated by floating-point arithmetic and
repeated force-calculation loops.

The RISC-V binary (\texttt{namd\_r.riscv}) was compiled with gem5 ROI
markers inserted around the main iteration loop in \texttt{spec\_namd.C}.
The markers delimit the Region of Interest using the gem5 \texttt{m5ops}
pseudo-instructions:

\begin{itemize}
    \item \texttt{m5\_work\_begin(0,0)} — inserted before the iteration loop
    \item \texttt{m5\_work\_end(0,0)}   — inserted after  the iteration loop
\end{itemize}

Only spills occurring between these two markers are counted toward the
reported statistics.

% =========================================================
\section{Simulation Setup}

\subsection{Input}

NAMD accepts a single molecular structure file (\texttt{apoa1.input})
shared across all dataset sizes. The dataset variant is controlled
solely by the \texttt{--iterations} argument. Table~\ref{tab:inputs}
summarises the three standard dataset sizes.

\begin{table}[H]
\centering
\caption{NAMD Input Dataset Sizes}
\label{tab:inputs}
\begin{tabular}{lcc}
\toprule
\textbf{Dataset} & \textbf{Iterations} & \textbf{Input File} \\
\midrule
test     & 1  & apoa1.input \\
train    & 8  & apoa1.input \\
refrate  & 65 & apoa1.input \\
\bottomrule
\end{tabular}
\end{table}

The \textbf{test} dataset (1 iteration) was selected for this run to
obtain a quick and verifiable result. The input file \texttt{apoa1.input}
is 7.8\,MB and represents the ApoA1 system with 92,224 atoms.

\subsection{gem5 Configuration}

\begin{itemize}
    \item \textbf{CPU model:}   TimingSimpleCPU
    \item \textbf{ISA:}         RISC-V RV64GC
    \item \textbf{Caches:}      enabled (default L1 I/D + L2)
    \item \textbf{Memory:}      4\,GB simulated physical RAM
    \item \textbf{Spill mode:}  summary only (\texttt{spill\_verbose=False})
\end{itemize}

The full gem5 invocation was:

\begin{verbatim}
./build/RISCV/gem5.opt
  --outdir=benchmarks/namd/m5out_test
  configs/deprecated/example/se.py
  --cmd=benchmarks/namd/namd_r.riscv
  --options="--input benchmarks/namd/apoa1.input
             --iterations 1
             --output benchmarks/namd/apoa1.test.output"
  --cpu-type=TimingSimpleCPU
  --caches
  --mem-size=4GB
\end{verbatim}

% =========================================================
\section{Simulation Results}

The simulation completed successfully with the output:
\texttt{"SUCCESSFUL COMPLETION — 1 ITERATIONS"}.

\subsection{Pre-ROI Context}

Before the ROI was entered, the binary executed setup and
initialization code. Table~\ref{tab:preroi} shows the global
counters at the moment \texttt{m5\_work\_begin} was triggered.

\begin{table}[H]
\centering
\caption{Global Counters at ROI Entry}
\label{tab:preroi}
\begin{tabular}{lr}
\toprule
\textbf{Metric} & \textbf{Value} \\
\midrule
Total instructions (pre-ROI) & 1,558,178,327 \\
Total spills detected (pre-ROI) & 125,215,610 \\
\bottomrule
\end{tabular}
\end{table}

\subsection{ROI Spill Statistics}

Table~\ref{tab:roi} presents the spill statistics collected exclusively
within the ROI (one NAMD iteration).

\begin{table}[H]
\centering
\caption{ROI Spill Statistics — 508.namd\_r (test, 1 iteration)}
\label{tab:roi}
\begin{tabular}{lr}
\toprule
\textbf{Metric} & \textbf{Value} \\
\midrule
ROI instructions         & 31,914,751,100 \\
ROI stores               &  1,945,521,360 \\
ROI loads                &  8,380,476,926 \\
\midrule
\textbf{ROI spills}      & \textbf{327,718,955} \\
\textbf{Spill rate}      & \textbf{1.0269\%} \\
\bottomrule
\end{tabular}
\end{table}

\subsection{Timing}

\begin{table}[H]
\centering
\caption{Simulation Timing}
\label{tab:timing}
\begin{tabular}{lr}
\toprule
\textbf{Metric} & \textbf{Value} \\
\midrule
Simulated time  & 68.73\,s \\
Host wall time  & 25,169\,s ($\approx$7\,hours) \\
\bottomrule
\end{tabular}
\end{table}

% =========================================================
\section{Analysis}

\subsection{Spill Rate Interpretation}

The measured spill rate of \textbf{1.03\%} means that approximately
1 in every 97 instructions inside the ROI is a register spill
(a stack load matching a prior stack store).

Within the ROI, the ratio of spills to total loads is:
\[
\frac{327{,}718{,}955}{8{,}380{,}476{,}926} \approx 3.91\%
\]
meaning roughly 1 in 25 load instructions is a spill reload.

\subsection{Store vs.\ Load Balance}

The load count (8.38B) is approximately 4.3$\times$ higher than the
store count (1.95B). This asymmetry is characteristic of
floating-point-heavy kernels: values are loaded many times for
repeated force calculations but written back to the stack
less frequently.

\subsection{Pre-ROI vs.\ ROI Spill Density}

Pre-ROI spills: $125{,}215{,}610$ over $1{,}558{,}178{,}327$ instructions
= \textbf{8.03\%} spill rate. This is notably higher than the ROI
rate (1.03\%), indicating that NAMD's initialization and data-loading
phase is more register-pressure-intensive than the core iteration loop.

\subsection{Scaling to Refrate}

The test dataset uses 1 iteration; the refrate dataset uses 65.
Assuming the ROI spill count scales linearly with iteration count,
the estimated refrate spill count would be:
\[
327{,}718{,}955 \times 65 \approx 21.3 \text{ billion spills}
\]
This projection should be verified against an actual refrate run.

% =========================================================
\section{Conclusion}

The \texttt{508.namd\_r} benchmark exhibits a moderate register spill
rate of 1.03\% within its core computation ROI on RISC-V RV64GC.
With 327 million spill events across 31.9 billion instructions in a
single iteration, the benchmark presents a meaningful but not dominant
spill workload. The pre-ROI initialization phase shows a significantly
higher spill density (8.03\%), suggesting that the molecular structure
loading phase is the primary source of register pressure in NAMD.

\end{document}
